%!TEX program = xelatex
\let\stop\empty

\documentclass[12pt,a4paper]{exam-zh}

\usepackage{xeCJK}
\setCJKmainfont{Noto Serif CJK SC}
\usepackage{geometry}
\usepackage{array}
\usepackage{tikz}
\usepackage{amsmath}
\usepackage{wrapfig}
\usepackage{adjustbox}
\usepackage{enumitem}
\usepackage{fancyhdr}
\usepackage{pgfplots}
\pgfplotsset{compat=1.18}
\usetikzlibrary{calc}

\usepackage{scalerel}
\def\mcirc{\mathbin{\scalerel*{\circ}{j}}}
\def\msquare{\mathord{\scalerel*{\Box}{gX}}}
\geometry{margin=2cm}

\AtBeginDocument{\zihao{4}}
% ----------- 调整 exam -----------


%region exam setup
\examsetup{
    % problem/label = \arabic*., 
    % 设置大题题号格式 
        question = {
        top-sep = 1em,      % 题头上方的间距（默认 0.25em）
        bottom-sep = 1em,   % 题号下方的间距（默认 0pt）
    },
    % 隐藏大题后的分值显示 
    problem/show-points = false,
    % 调整小题列表的缩进和间距
    list/step-name = {}, 
    solution = {
    show-solution = show-move, % 答案自动移动到卷末
    label-content = {answer},  % 答案区的标题文字
    show-qed = true            % 在解答末尾显示结束符号 ◻ 
  },
    page = {
    size = a4paper,           % 设定纸张大小 
    foot-content = { - ; - }  % 页脚显示为 -1-, -2- 格式 
  }
}
\setlist[enumerate,1]{
    label=(\alpha*), 
    itemsep=0.5em, 
    leftmargin=2em, 
    after=\hfill % 配合您之前要求的百分比靠右样式
}
\newlist{questionparts}{enumerate}{1}
\setlist[questionparts]{
    label=(\alph*),
    itemsep=0.5em,
    leftmargin=2em,
    rightmargin=4em
}
\newcommand{\qscore}[1]{%
    \unskip\nobreak\hfill
    \makebox[0pt][l]{\makebox[4em][r]{(#1\%)}}%
}
\newcommand{\solutionlayout}{%
    \setlength{\abovedisplayskip}{0.6em}%
    \setlength{\belowdisplayskip}{0.6em}%
    \setlength{\abovedisplayshortskip}{0.35em}%
    \setlength{\belowdisplayshortskip}{0.35em}%
    \setlength{\jot}{0.35em}%
}
%endregion




\begin{document}

%region 范围：选择性必修第三册·成对数据的统计分析 | 考点：样本相关系数的计算 | 难度：基础
\begin{question}
    已知关于 $X$ 与 $Y$ 的 $50$ 组观测数据 $(x_i,y_i)\ (i=1,2,\cdots,50)$，其中
    $        \displaystyle
        \sum\limits_{i=1}^{50}x_i=123,\quad \sum\limits_{i=1}^{50}x_i^2=12321,\quad
        \sum\limits_{i=1}^{50}y_i=243,\quad \sum\limits_{i=1}^{50}y_i^2=67854,\quad
        \sum\limits_{i=1}^{50}x_iy_i=23436,$

    求 $X$ 与 $Y$ 的样本相关系数 $r$（结果保留两位小数）。
\end{question}

\begin{solution}
    求 $X$ 与 $Y$ 的样本相关系数 $r$。\\
    利用相关系数公式，先求均值，再求三个离差平方和。\\
    \begin{align*}
        r &= \frac{\displaystyle\sum\limits_{i=1}^{n}x_iy_i-n\bar{x}\,\bar{y}}
        {\sqrt{\left(\displaystyle\sum\limits_{i=1}^{n}x_i^2-n\bar{x}^2\right)\left(\displaystyle\sum\limits_{i=1}^{n}y_i^2-n\bar{y}^2\right)}}
    \end{align*}
    $n=50$，则
    \begin{align*}
        \bar{x} &= \frac{123}{50}=2.46, \qquad \bar{y}=\frac{243}{50}=4.86 \\
        \sum\limits_{i=1}^{50}x_i^2-n\bar{x}^2 &= 12321-50\times 2.46^2=12018.42 \\
        \sum\limits_{i=1}^{50}y_i^2-n\bar{y}^2 &= 67854-50\times 4.86^2=66673.02 \\
        \sum\limits_{i=1}^{50}x_iy_i-n\bar{x}\,\bar{y} &= 23436-50\times 2.46\times 4.86=22838.22 \\
        r &= \frac{22838.22}{\sqrt{12018.42\times 66673.02}} \\
        &\approx \frac{22838.22}{28307.3} \\
        &\approx 0.81
    \end{align*}
\end{solution}
%endregion

%region 范围：选择性必修第三册·成对数据的统计分析 | 考点：散点图、样本相关系数的计算与相关程度判断 | 难度：基础
\begin{question}
    为研究人体心脏的收缩压 $Y$（单位：mmHg）与年龄 $X$（单位：岁）之间的关系，某医学院随机收集了 $13$ 名男子的数据，如下表所示：
    \begin{center}
        \small
        \begin{tabular}{c|*{13}{c}}
            \hline
            年龄 $x_i$（岁）  & 51  & 22  & 23  & 31  & 33  & 49  & 58  & 53  & 44  & 55  & 42  & 45  & 25  \\
            \hline
            收缩压 $y_i$（mmHg） & 130 & 141 & 124 & 126 & 117 & 135 & 143 & 138 & 132 & 143 & 133 & 155 & 147 \\
            \hline
        \end{tabular}
    \end{center}

    \begin{questionparts}
        \item 作收缩压与年龄的散点图；

        \item 求收缩压与年龄的样本相关系数 $r$（结果保留两位小数），并判断两者的相关程度。
    \end{questionparts}
\end{question}

\begin{solution}
    \begin{enumerate}[label=(\alph*)]
        \item 作 $13$ 名男子的年龄与收缩压的散点图。\\
        以年龄为横坐标、收缩压为纵坐标，描出各点 $(x_i,y_i)$。
        \begin{center}
            \begin{tikzpicture}[x=0.2cm, y=0.1cm, >=stealth]
                % 坐标轴
                \draw[->] (20,110) -- (63,110) node[right] {年龄/岁};
                \draw[->] (20,110) -- (20,164) node[above] {收缩压/mmHg};
                % 刻度
                \foreach \x in {20,30,40,50,60}{
                    \draw (\x,110) -- (\x,112);
                    \node[below] at (\x,110) {\small \x};
                }
                \foreach \y in {110,120,130,140,150,160}{
                    \draw (20,\y) -- (21,\y);
                    \node[left] at (20,\y) {\small \y};
                }
                % 数据点
                \foreach \x/\y in {51/130,22/141,23/124,31/126,33/117,49/135,58/143,53/138,44/132,55/143,42/133,45/155,25/147}{
                    \fill (\x,\y) circle (1.8pt);
                }
            \end{tikzpicture}
        \end{center}

        \item 求收缩压与年龄的样本相关系数 $r$，并判断相关程度。\\
        利用 $\displaystyle r=\frac{n\sum\limits_{i=1}^{n}x_iy_i-\sum\limits_{i=1}^{n}x_i\sum\limits_{i=1}^{n}y_i}{\sqrt{\left[n\sum\limits_{i=1}^{n}x_i^2-\left(\sum\limits_{i=1}^{n}x_i\right)^2\right]\left[n\sum\limits_{i=1}^{n}y_i^2-\left(\sum\limits_{i=1}^{n}y_i\right)^2\right]}}$ 计算，再依 $|r|$ 的大小判断。\\
        $n=13$，则
        \begin{align*}
            \sum\limits_{i=1}^{13}x_i &= 531, \quad \sum\limits_{i=1}^{13}y_i=1764 \\
            \sum\limits_{i=1}^{13}x_i^2 &= 23613, \quad \sum\limits_{i=1}^{13}y_i^2=240636, \quad \sum\limits_{i=1}^{13}x_iy_i=72483 \\
            r &= \frac{13\times 72483-531\times 1764}{\sqrt{\left(13\times 23613-531^2\right)\left(13\times 240636-1764^2\right)}} \\
            &= \frac{5595}{\sqrt{25008\times 16572}} \\
            &\approx \frac{5595}{20357.6} \\
            &\approx 0.27
        \end{align*}
        $0<r<0.3$，故收缩压与年龄呈正相关，但线性相关程度很弱。
    \end{enumerate}
\end{solution}
%endregion

%region 范围：选择性必修第三册·成对数据的统计分析 | 考点：样本相关系数、回归直线方程及预测 | 难度：基础
\begin{question}
    某班 $16$ 名学生参加一次音乐与体育的评估，音乐分数 $X$ 的满分为 $60$ 分，体育分数 $Y$ 的满分为 $40$ 分。成绩简报如下：
    
      $  \sum\limits_{i=1}^{16}x_i=682,\quad \sum\limits_{i=1}^{16}x_i^2=29982,\quad
        \sum\limits_{i=1}^{16}y_i=407,\quad \sum\limits_{i=1}^{16}y_i^2=10815,\quad
        \sum\limits_{i=1}^{16}x_iy_i=17060.$
    

    \begin{questionparts}
        \item 求音乐与体育分数的样本相关系数 $r$（结果保留两位小数），并判断其相关程度；

        \item 求 $Y$ 对 $X$ 的回归直线方程（系数保留四位小数，常数项保留两位小数）；

        \item 若小芬的音乐得分为 $55$ 分，预测她的体育得分（保留一位小数）。
    \end{questionparts}
\end{question}

\begin{solution}
    \begin{enumerate}[label=(\alph*)]
        \item 求音乐与体育分数的样本相关系数 $r$，并判断相关程度。\\
        用 $r=\dfrac{n\sum x_iy_i-\sum x_i\sum y_i}{\sqrt{\left[n\sum x_i^2-\left(\sum x_i\right)^2\right]\left[n\sum y_i^2-\left(\sum y_i\right)^2\right]}}$ 计算，再依 $|r|$ 判断。
        \begin{align*}
            n\sum\limits_{i=1}^{16}x_iy_i-\sum\limits_{i=1}^{16}x_i\sum\limits_{i=1}^{16}y_i &= 16\times 17060-682\times 407=-4614 \\
            n\sum\limits_{i=1}^{16}x_i^2-\left(\sum\limits_{i=1}^{16}x_i\right)^2 &= 16\times 29982-682^2=14588 \\
            n\sum\limits_{i=1}^{16}y_i^2-\left(\sum\limits_{i=1}^{16}y_i\right)^2 &= 16\times 10815-407^2=7391 \\
            r &= \frac{-4614}{\sqrt{14588\times 7391}} \\
            &\approx \frac{-4614}{10383.6} \\
            &\approx -0.44
        \end{align*}
        $0.3\le|r|<0.75$，故音乐与体育分数呈中度负相关。

        \item 求 $Y$ 对 $X$ 的回归直线方程 $\hat{y}=bx+a$。\\
        用 $b=\dfrac{n\sum x_iy_i-\sum x_i\sum y_i}{n\sum x_i^2-\left(\sum x_i\right)^2}$，$a=\bar{y}-b\bar{x}$。
        \begin{align*}
            \bar{x} &= \frac{682}{16}=42.625, \qquad \bar{y}=\frac{407}{16}=25.4375 \\
            b &= \frac{-4614}{14588}\approx -0.3163 \\
            a &= 25.4375-(-0.3163)\times 42.625\approx 38.92 \\
            \hat{y} &= -0.3163x+38.92
        \end{align*}

        \item 当 $x=55$ 时，预测体育得分 $\hat{y}$。\\
        将 $x=55$ 代入回归直线方程。
        \begin{align*}
            \hat{y} &= -0.3163\times 55+38.92 \\
            &\approx 21.5
        \end{align*}
    \end{enumerate}
\end{solution}
%endregion

%region 范围：选择性必修第三册·成对数据的统计分析 | 考点：样本相关系数、回归直线方程及预测 | 难度：基础
\begin{question}
    下表为 $10$ 个同类企业的固定资产价值 $X$ 与总产值 $Y$ 的数据（单位：万令吉）：
    \begin{center}
        \small
        \begin{tabular}{c|*{10}{c}}
            \hline
            企业编号 & 1 & 2 & 3 & 4 & 5 & 6 & 7 & 8 & 9 & 10 \\
            \hline
            固定资产价值 $x_i$ & 200 & 314 & 318 & 409 & 415 & 502 & 910 & 1022 & 1210 & 1225 \\
            \hline
            总产值 $y_i$ & 638 & 605 & 524 & 815 & 913 & 928 & 1019 & 1219 & 1516 & 1624 \\
            \hline
        \end{tabular}
    \end{center}

    \begin{questionparts}
        \item 求固定资产价值与总产值的样本相关系数 $r$（结果保留两位小数），并判断其相关程度；

        \item 求 $Y$ 对 $X$ 的回归直线方程（系数保留四位小数，常数项保留两位小数）；

        \item 若一家同类企业的固定资产价值为 $800$ 万令吉，预测其总产值（保留一位小数）。
    \end{questionparts}
\end{question}

\begin{solution}
    \begin{enumerate}[label=(\alph*)]
        \item 求固定资产价值与总产值的样本相关系数 $r$，并判断相关程度。\\
        先求各项和，再代入 $r=\dfrac{n\sum x_iy_i-\sum x_i\sum y_i}{\sqrt{\left[n\sum x_i^2-\left(\sum x_i\right)^2\right]\left[n\sum y_i^2-\left(\sum y_i\right)^2\right]}}$。\\
        $n=10$，则
        \begin{align*}
            \sum\limits_{i=1}^{10}x_i &= 6525, \quad \sum\limits_{i=1}^{10}y_i=9801 \\
            \sum\limits_{i=1}^{10}x_i^2 &= 5668539, \quad \sum\limits_{i=1}^{10}y_i^2=10866577, \quad \sum\limits_{i=1}^{10}x_iy_i=7659156 \\
            n\sum\limits_{i=1}^{10}x_iy_i-\sum\limits_{i=1}^{10}x_i\sum\limits_{i=1}^{10}y_i &= 10\times 7659156-6525\times 9801=12640035 \\
            n\sum\limits_{i=1}^{10}x_i^2-\left(\sum\limits_{i=1}^{10}x_i\right)^2 &= 10\times 5668539-6525^2=14109765 \\
            n\sum\limits_{i=1}^{10}y_i^2-\left(\sum\limits_{i=1}^{10}y_i\right)^2 &= 10\times 10866577-9801^2=12606169 \\
            r &= \frac{12640035}{\sqrt{14109765\times 12606169}} \\
            &\approx \frac{12640035}{13336796} \\
            &\approx 0.95
        \end{align*}
        $|r|\ge 0.75$，故固定资产价值与总产值呈高度正相关。

        \item 求 $Y$ 对 $X$ 的回归直线方程 $\hat{y}=bx+a$。\\
        用 $b=\dfrac{n\sum x_iy_i-\sum x_i\sum y_i}{n\sum x_i^2-\left(\sum x_i\right)^2}$，$a=\bar{y}-b\bar{x}$。
        \begin{align*}
            \bar{x} &= \frac{6525}{10}=652.5, \qquad \bar{y}=\frac{9801}{10}=980.1 \\
            b &= \frac{12640035}{14109765}\approx 0.8958 \\
            a &= 980.1-0.8958\times 652.5\approx 395.57 \\
            \hat{y} &= 0.8958x+395.57
        \end{align*}

        \item 当 $x=800$ 时，预测总产值 $\hat{y}$。\\
        将 $x=800$ 代入回归直线方程。
        \begin{align*}
            \hat{y} &= 0.8958\times 800+395.57 \\
            &\approx 1112.2
        \end{align*}
        即预测总产值约为 $1112.2$ 万令吉。
    \end{enumerate}
\end{solution}
%endregion

\end{document}